\documentclass[../../main-detail.tex]{subfiles}
\begin{document}
\chapter{表現と虚構物 -- kika以降}\label{chap:kika}
本章では，\kotr{kika}{準抽象物}以下の分類と，表現・実現・虚構物・形式的構造の関係を定める．
答える問いは4つである：小説や楽譜は何でできているか（表現の構成），同じ小説が複数の冊子にあるとは
どういうことか（実現），ホームズは何であって何でないか（虚構物），文字列や数はどこに置くか（抽象物との境界）．

\begin{remark}[2026-08-29改訂]
旧草案（2026-08-01，K1--K8）から次を変更した．
\begin{itemize}
    \item \ko{kilemkto}（表現実現物）を廃止し，\ko{saden}の値域を具体物全般に広げた．
    \item 形態\ko{kilemhaisk}を\ko{kika}から\ko{stcilkant}へ移した．文字列はずっとあるもの，俳句はある時点で生み出されたものとして区別する．
    \item 符号化規約$E_{\mathrm{code}}$を個体として立てた（未造語）．
    \item 虚構物の層Kをタイプ個体（上位分類の章「種を立てる基準」）と同一視した．
    \item 実現連鎖に推移性を課した．
    \item 旧\ko{takht}「$x_2$は$x_1$という命題を主張している」は，符号化・aboutness・発話の力を一関係に混ぜていたため廃止する．
\end{itemize}
K1（表現の複合体化），K2（符号化とaboutnessの分離），K3（\ko{kilem}再定義），K4（虚構性を種にしない），K5（由来と参照の分離）は採択した．
K6（\ko{kika}作業定義）は暫定のまま，K7（語義分割）は方針として維持，K8（語彙テスト）は未実施である．
\end{remark}

\section{対象と凡例}

\begin{figure}[H]
  \centering
  \begin{forest}
    dir tree
    [(root), draw=none
      [\tn{toika}{具体物}
        [既存の木（\ko{nato}・\ko{aicis}・\ko{astav}），実現物はここ, draw=none, edge={dashed}]
      ]
      [\tn{kika}{準抽象物}
        [{\tn{kilem}{表現} $R=\langle F_{\mathrm{form}},C_{\mathrm{cont}},E_{\mathrm{code}}\rangle$}]
        [\tn{klaleme}{内容} $C_{\mathrm{cont}}$]
        [(符号化規約 $E_{\mathrm{code}}$)：言語・暗号鍵・記法]
        [(作品)：表現を束ねる歴史的タイプ個体（暫定）]
        [(キャラクター役)：層K＝タイプ個体]
        [(記述)：制度を構成する規約の個体．法人]
      ]
      [\tn{loka}{抽象物}
        [\tn{latent}{量}, draw=none, edge={dashed}]
      ]
      [\tn{stcilkant}{形式的構造}
        [\tn{kilemhaisk}{形態} $F_{\mathrm{form}}$：文字列・旋律・線画]
      ]
    ]
  \end{forest}
  \caption{表現と虚構物に関わる分類．括弧内は未造語．}
\end{figure}

各クラスについて，トークンが何であるか（凡例）と，何で個体化するか（同一性基準）を先に宣言する．
\begin{itemize}
    \item \emph{\ko{kika}}：トークンは，制作・制定・解釈の歴史によって個体化され，媒体・心的状態・共同体のいずれかに一般的に依存する反復可能な対象．
    それ自体の同一性では空間に局在せず，同一性が形式体系の内部で完結しない（作業定義，暫定）．
    \item \emph{\ko{kilem}（表現）}：トークンは形態$F_{\mathrm{form}}$・内容$C_{\mathrm{cont}}$・符号化規約$E_{\mathrm{code}}$の組で，制作史を持つもの．
    \ko{kilem}は開ハブであり，形態・内容・符号化規約への三つの原始射影を持つ．三射影の一致は同一性の十分条件ではなく，制作史の異なる表現トークンを区別できる．同一性は原始的様式で運用する（2026-09-03．旧「三成分の組＋制作史という剰余同一性」は開／閉ハブの二分に乗らないため改めた）．
    \ko{canon}が担い手・値への射影を持つ開ハブであるのと同じ型である（量・性質モジュールの章）．
    \item \emph{\ko{kilemhaisk}（形態）}：トークンは文字列・旋律・線画・配置パタンなど，内容を担いうる反復可能な形態．
    \ko{stcilkant}に置く．同一性は許容変形の同値関係に依存し，フォント差を無視する文字列と字体が本質のロゴでは別の同値関係を使う．
    どの同値関係かは語彙ごとに記録する．
    \item \emph{\ko{klaleme}（内容）}：トークンは命題，物語内容，規則，設計内容など，形態によって符号化されるもの．
    同一性は論理的同値では尽きず，翻訳・編曲・改訂をどこまで同じ内容と数えるかは内容種別ごとの規約による．
    \item \emph{符号化規約$E_{\mathrm{code}}$}（未造語）：トークンは言語（ラング），暗号鍵，記法慣習，文字コードなど．
    外延は$F\times C$の部分集合（形態から内容への部分関数のグラフ）であり，制定史によって個体化される．
    \item \emph{作品}（未造語，暫定）：制作・公表・継承の歴史で個体化され，複数の表現を束ねるタイプ個体．\ko{kilem}と別に立て，表現との関係はExpressionOf（\ref{sec:fiction}節．2026-09-06）．
    \item \emph{キャラクター役}（未造語）：人物・設定などのタイプ個体（\ref{sec:fiction}節）．作品に依存せず立つ（初音ミク）．
    \item \emph{〈記述〉}（未造語，暫定）：会社・国家・委員会などの制度を構成する規約の個体．設立・制定・承継の歴史で個体化され，定款や会則の改訂を経ても同じ個体として続く（作品と同じく内容の同一性では個体化しない）．構成される組織（\ko{takto}）が存在しない時期にも存続する．法人は法的地位を持つ〈記述〉である．組織との関係は見做しで与える（持続物の章「組織と法人」．2026-09-28）．
    \item \emph{実現物}：クラスを立てない．冊子・録音媒体は\ko{toika}の既存の木に，朗読・演奏・表示は\ko{astav}に普通に分類し，
    表現との関係は\ko{saden}だけが担う．
\end{itemize}

\section{表現の構成}
表現は形態と内容を符号化規約で結んだ複合体$R=\langle F_{\mathrm{form}},C_{\mathrm{cont}},E_{\mathrm{code}}\rangle$である
\citep{yamato2022,mizoguchi-borgo2022representation,mizoguchi-borgo2025representation}．
ソシュールの二項（シニフィアン$F_{\mathrm{form}}$，シニフィエ$C_{\mathrm{cont}}$）に対して，両者を結ぶ体系を第三項として個体化する
（Hjelmslevの記号関数，Ecoのコードに近い）．

\paragraph{関係} 3成分への関係はどれも$R$からの射影で，関数的である．
\rele{meelja}{kilemhaisk}{kilem}{形態-表現関係，$x_2$の形態は$x_1$である}
\rele{miz}{klaleme}{kilem}{内容-表現関係，$x_2$の内容は$x_1$である}
規約参照（$E_{\mathrm{code}}$-表現関係）は未造語である．
$E_{\mathrm{code}}$は形態と内容の間のチャートであり，単位関数cm・m（量$\to$数値）やカレンダー（時$\to$日付）と同じ族に属する．
同じ形態が慣習や鍵の違いで別の内容を符号化するとき，$F_{\mathrm{form}}$と$C_{\mathrm{cont}}$の組が同じでも$E_{\mathrm{code}}$が違えば別の表現である．

\paragraph{言語の立ち位置} 「言語」は二つに割れる．規約の束としての言語（ラング）は$E_{\mathrm{code}}$側の\ko{kika}トークン，
その言語で書ける文字列全体は$F_{\mathrm{form}}$側の類（\ko{stcilkant}）である．旧版は両者を「言語は\ko{kilemhaisk}に含まれる」で混ぜていた．
コノメノ自身の表現については，フレームの語彙割り当て$\rho$（形式文法章）がコノメノ表現の$E_{\mathrm{code}}$に当たる．

\paragraph{符号化・aboutness・力を分ける} $F_{\mathrm{form}}$が$C_{\mathrm{cont}}$を符号化することと，$C_{\mathrm{cont}}$が何かについてのものであること，
$C_{\mathrm{cont}}$が主張・命令として発せられることは三つの別の関係である．
本章は符号化のみを扱い，aboutnessは未造語の関係として内容から対象（世界内個体，表現，状況）へ引き，
発話の力は談話モジュールの発話規則に置く．DnSのDescription--Situation--satisfies\citep{gangemi-mika2003}はaboutness側の比較対象である．

\section{実現}
\rele{saden}{kilem}{haltoika}{実現関係，$x_2$は$x_1$を実現する}
\ko{saden}の値域は\ko{haltoika}（広義具体物）$=$\ko{toika}$\cup$\ko{kika}である（2026-09-03．表現を実現する表現——版・刷——を値域に含めるため．\ko{haltoika}は生起物の章の\ko{stae}の定義域と同じ「時間に依存する個体」の類で，辞書v5には未登録）．朗読や演奏のような生起物による実現も同じ関係で扱う
（\ko{astav}は\ko{toika}の下位なので値域内であり，担い手で代表させる必要はない）．
何かを表現の実現と見なすかは評価者・記述に相対的なので（岩層を記録として読む），実現物のクラスは立てず，
相対性は\ko{saden}の真理条件には組み込まない（2026-09-07裁定）：\ko{saden}自体は実現関係であり，ある評価者が実現と認めることは実現命題$K[\mathrm{saden}(R,x)]$を内容とする判断（内容対象方式）として表し，\ko{saden}の事実とは区別する．「記録として読める」は，解読規約・内容・形態を備えた表現$R$を指定規約の下で$x$から解読できるという能力の記述であり，実際の実現判断とも区別する．判断・能力のどちらからも無条件の$\mathrm{saden}(R,x)$は導かない（埋め込みを除去する推論は認めない）．評価者ごとの関係へ\ko{saden}を再定義すると異なる基準の実現連鎖を合成する条件まで変わるので採らない．「$X$として」の一般論は量・性質モジュールの章．

\paragraph{推移性} 実現は連鎖する：原文の表現$\to$版$\to$刷$\to$冊子，ファイル$\to$表示．
中間の駅を立てるかは恣意的なので，\ko{saden}に推移性を課し，原文$\to$冊子と原文$\to$版$\to$刷$\to$冊子を両立させる（連鎖の上端は\ko{kilem}であり作品ではない．作品から実現物への関係はExpressionOfと\ko{saden}の合成で定義する．\ref{sec:fiction}節，2026-09-06）．推移律は中間項が\ko{kilem}（\ko{kika}側）である連鎖に限る．値域を\ko{haltoika}に広げたのはそのためである．楽譜$\to$演奏$\to$録音のように中間項が\ko{toika}（演奏は\ko{astav}）である連鎖には推移律を適用せず，録音は演奏の記録であって楽譜の実現とは言わない（2026-09-03）．
これは方法論の「粒度中立性」の実例で，中間の駅をどの粒度で個体化しても上流の実現事実は変わらない．
版・刷などの名前付き切片は保存された検索として扱い，木のノードにしない．
実現の連鎖では，実現物そのものが形態になるのではなく，実現物から許容変形で同一視された形態トークンを抽出し，それを別の\ko{kilem}の\ko{meelja}射影へ渡す（次節）．この抽出関係は未造語である．

\section{再帰}\label{sec:rep-recursion}
\ko{canon}の再帰は担い手スロットが\ko{canon}を受け入れる一チャネルだった．表現には三チャネルある．
\begin{enumerate}
    \item[(a)] \emph{内容スロット}：$C_{\mathrm{cont}}$に表現や形態が入る．引用，報告，暗号の入れ子，翻案（内容が別の表現に由来する）．
    \item[(b)] \emph{実現の連鎖}：実現物がまた上位の形態トークンになる（前節）．
    \item[(c)] \emph{aboutness}：内容が表現についてのものである．批評，注釈，書誌．
\end{enumerate}
三チャネルとも，レベルごとに新語を立てず同じ関係の再適用で書く（\ko{nilof}/\ko{noclof}を高階の\ko{ansa}/\ko{noccea}に還元したのと同じ処方）．

\begin{figure}[H]
  \centering
  \begin{tikzpicture}[
      n/.style={draw, rounded corners=2pt, align=center, font=\small, inner sep=4pt},
      r/.style={-{Stealth}, font=\footnotesize},
      rec/.style={-{Stealth}, blue!60!black, font=\footnotesize, text=blue!60!black}
    ]
    \node[n] (R)   at (0,0)       {\ko{kilem} $R=\langle F_{\mathrm{form}},C_{\mathrm{cont}},E_{\mathrm{code}}\rangle$};
    \node[n] (F)   at (-4.4,2.2)  {\ko{kilemhaisk} $F_{\mathrm{form}}$\\（\ko{stcilkant}）};
    \node[n] (C)   at (0,3.0)     {\ko{klaleme} $C_{\mathrm{cont}}$};
    \node[n] (E)   at (4.4,2.2)   {（符号化規約 $E_{\mathrm{code}}$）};
    \node[n] (T)   at (-4.4,-2.2) {実現物（\ko{toika}）};
    \node[n] (K)   at (4.6,-2.2)  {（キャラクター役）\\層K＝タイプ個体};
    \node[n] (Wld) at (1.4,-4.2)  {層W個体};
    \node[n] (Cls) at (-2.6,-4.2) {充足世界クラス};
    \node[n] (Act) at (-6.4,0.2)  {制作行為（\ko{astav}）};
    \draw[r] (F) -- node[left=1pt]{meelja} (R);
    \draw[r] (C) -- node[right=1pt]{miz} (R);
    \draw[r] (E) -- node[right=1pt]{規約参照} (R);
    \draw[r] (R) -- node[above=1pt, sloped]{saden（推移的）} (T);
    \draw[r] (Act) -- node[above=1pt, sloped]{歴史的個体化} (R);
    \draw[r] (K) -- node[right=2pt]{対応（ハブ経由）} (Wld);
    \draw[r] (C) .. controls (-2.2,0.8) and (-2.4,-2.6) .. node[left=1pt, pos=0.8]{$K[\phi_d]$} (Cls);
    \draw[r] (C) .. controls (2.0,0.6) and (2.2,-2.6) .. node[right=1pt, pos=0.85]{about} (Wld);
    \draw[rec] (R) .. controls (1.8,1.4) .. node[right=2pt]{(a) 引用・翻案} (C);
    \draw[rec] (T) .. controls (-6.2,-0.8) and (-6.4,1.2) .. node[left=2pt]{(b) 連鎖} (F);
    \draw[rec] (C) .. controls (3.0,3.4) and (2.6,1.2) .. node[right=2pt, pos=0.15]{(c) 批評} (R);
  \end{tikzpicture}
  \caption{表現と虚構物の主な関係．青は再帰チャネル．}
\end{figure}

\paragraph{形式意味論との同定（暫定）} チャネル(a)は形式意味論章で既に実装されている見込みが高い：
形態\ko{kilemhaisk}は構文オブジェクト$\mathrm{Rep}(\alpha)$（urelementへの埋め込み），引用の再帰は$E^\omega$と冠詞$M$，
自己言及表現の存在は終余代数的符号化の採択に対応する．表現の存在論は新しい理論ではなく，意味論側の既存装置の
オントロジー的読み替えとして書けるはずである（同定表は未決事項）．

\section{虚構物}\label{sec:fiction}
虚構物の記述では，世界内での対象，歴史的に個体化された内容・文化的人工物，形式的構造を区別し，それぞれ層W，層K，層Lと呼ぶ（2026-09-06改訂，暫定．旧稿は「三つの層に分かれる」としていたが，三層は存在者を排他的に分割する上位分類ではなく，記述に現れる対象を区別するための図式である）．
\begin{itemize}
    \item \emph{層W（世界内個体）}：世界$w$の中で探偵として活動するホームズ．当該世界の既存の分類に従う．人物・道具は\ko{nato}，出来事は\ko{astav}，個別の性質・能力は\ko{aicis}であり，作中で制作された表現はその世界でも\ko{kika}である．「層Wの対象$\subseteq$\ko{toika}」という一律の制約は課さない．
    \item \emph{層K（内容・文化的人工物）}：作品，表現，内容，設定，キャラクター役．制作・公表・改訂の歴史を持つ\ko{kika}．このうちキャラクター役と設定はタイプ個体であり，層Kの全対象がタイプ個体なのではない．
    \item \emph{層L（形式的構造）}：グラフ，群，力学系．\ko{stcilkant}．
\end{itemize}
虚構性は，ある世界における実例の欠如，内容との対応，話者の視点に関わる横断的な身分であり，種ではない．したがって\ko{kika}に「虚構物」というノードを設けない．どの対象にも三層すべての相手が存在するとは限らない．

\paragraph{キャラクター役は作品から独立したタイプ個体である} 層Kのキャラクター役（GOLEMのCharacter-Stoff，メタフィクション理論の「仕様」に当たる）は，上位分類の章「種を立てる基準」のタイプ個体であり，制作・継承・解釈の歴史に基づいて計数される．種を立てる診断がそのまま当たる：種水準の述定がある（「ホームズは1887年に作られた」），空外延で区別が要る（現実世界に実例のないホームズとポワロは別），「このキャラクター」への照応・計数がある．
役は作品に依存しない．初音ミクのように登場する作品を持たずに立つ役があるので，（登場）$\subseteq$作品$\times$役は依存関係ではなく登場の記録である．役は語り手・主人公などのロールタイプ（\ko{liflom}）とも別で，一つの役が複数の物語上のロールを担い，同じロールを異なる役が担いうる．

\paragraph{対応はハブで表す} 役と世界内個体の対応は，対応づける根拠を明示した引き下げハブ$H$で表す（2026-09-06採択，暫定）．ハブは役，対象への射影を持ち，必要に応じて根拠内容，物語行為，解釈行為への二項関係を外側に付加する．ハブの同一性は射影の一致では定めず，開ハブとして扱う．
\[
\begin{aligned}
&\text{役}\subseteq H\times T,\quad \text{対象}\subseteq H\times \mathcal M,\quad \text{根拠内容}\subseteq H\times\ko{klaleme},\quad \text{物語行為}\subseteq H\times\ko{astav},\quad \text{解釈行為}\subseteq H\times\ko{astav};\\
&\text{対応}(k,x)\ :\Longleftrightarrow\ \exists h\,[\,H(h)\land\text{役}(h,k)\land\text{対象}(h,x)\,].
\end{aligned}
\]
対応者関係はこのハブから誘導される二項関係で，単射性・関数性・全域性を課さない（多対多）．ハブを経由させる理由は設計規範の立証責任に答える：対応づけが$(a,x,d_1),(a,y,d_2),(b,x,d_2)$の三つだけのとき，役—対象，役—内容，対象—内容の裸の二項射影からは存在しない$(a,x,d_2)$も復元されてしまうので，引き下げなしには書けない．閉じた三成分束は不要である．
根拠を省いた対応は検索には使えるが，述定を移す範囲を決めない．「ホームズは勇敢である」はインスタンス水準の述定で，根拠内容を限定した対応先へ分配される（原作で勇敢，翻案で臆病ならハブが二つあり，役への直接述定にすると矛盾する）．「1887年に作られた」は種水準で層Kに留まる．対応先が存在しないことから全称的な作中述定を肯定してはならない．
ハブは対象を一意に生成するコンストラクタではない．対象が未確定の場合は，対応条件を満たす候補をフレームと談話モジュールの談話指示子で保持する．同じ役が現れることだけから，異なる叙述での対象の同一性を導かない．翻案由来$\subseteq$役$\times$役（および作品$\times$作品）は層K内の関係であり，対応者関係と区別する．翻案関係だけから同一性や性質の保存を導かない．

\paragraph{作品世界と背景の持ち込み} 作品について語る際には，設定内容，内容を成立させる状況，内容が真となる世界のクラスを区別する．設定内容は\ko{klaleme}の個体であり，世界クラスとの同一性では個体化しない（世界クラスが一致する異なる内容があり，世界クラスが空でも内容個体の区別は失われない）．
内容$d$から世界クラスを得るには，どの叙述を誰の主張として読み，どの事項を当該解釈の確定事項とするかを定め，その真理条件を表す文$\phi_d$に冠詞$K$を適用する：$W_d := K[\phi_d]$．冠詞$K$は内容個体を直接入力する演算子ではない．冠詞$C$は$\phi_d$の正確検証者を，冠詞$E$は既存の定義による選択肢を取り出す．解釈の違いが自動的に$E$の選択肢になるわけではない．
作中帰結は，確定事項と，その解釈で採用する背景を合わせて判定する（2026-09-06採択，暫定）．$K[\phi_d]$だけでは「作中でも地球は丸い」が出ない．背景候補$\Delta$（有限個の内容個体，各$b$の真理条件文を$\beta_b$）を記録し，
\[
U_J := K[\phi_d]\cap\bigcap_{b\in J}K[\beta_b],\quad \mathrm{Adm}:=\{J\subseteq\Delta\mid U_J\neq\emptyset\},\quad \mathrm{Max}:=\text{$\mathrm{Adm}$の包含極大元},\quad
\text{作中帰結}(d,\psi)\ :\Longleftrightarrow\ \mathrm{Max}\neq\emptyset\land\forall J\in\mathrm{Max}\,[\,U_J\subseteq K[\psi]\,]
\]
とする．すなわち確定事項と両立する包含極大の背景候補集合をすべて残し，そのすべてで成立する文だけを確定的な帰結とする．後の叙述で平面世界であることが確定すれば，「地球は丸い」を含む$J$が落ちる．等優先度の背景が競合すれば一方を勝手に選ばない．背景の出所と採用範囲は内容・行為への関係で記録し，現実世界の全事実や辞書の現実の分類辺を一括して背景にしてはならない．この方針は背景候補とともに明示し，一般的な物語理解の唯一の規則とはしない．Lewisの世界類似性による分析（Truth in Fiction, 1978）とは同値でなく，類似性が必要な用途でも比較基準を個体化し比較関係を引き下げれば書けるので，フレームに選択関数を新設しない（設計規範）．
内容が不整合で世界クラスが空になる場合，この方法による作中帰結を認めない．局所的に帰属された叙述は表現・内容・主張行為として保持できるが，不整合な作品全体についての非爆発的推論は現行の体系が提供しない．

\paragraph{非実在} 「キャラクターが実在しない」は，その役について認められる現実の指示先が存在しないという主張である：
\[
\text{現実指示}\subseteq T\times\ko{toika},\qquad \text{非実在}(k)\ :\Longleftrightarrow\ \lnot\exists x\,\text{現実指示}(k,x).
\]
作品全体の内容が現実世界で成立しないこと（$w_0\notin W_d$）からは導かない（実在人物が登場する架空の逸話が反例）．現実の指示先，創作の着想源，似た特徴を持つ人物は別の関係であり，現実指示を特徴一致で定義しない（旧稿の「充足世界クラスに現実世界が入らない」という説明は置換した）．

\paragraph{語り手・視点・作中作} 作者，語り手，視点の担い手を区別する．作者は表現の制作に参与し（\ko{ja}），語り手は叙述行為を担い，視点の担い手は知覚・知識・経験の帰属先である．同じ個体が複数を担っても関係は同一視しない．語り手がある内容を述べたことから，その内容が作中で成立することを導かない．人格的な語り手を同定しない叙述について，語り手ロールのプレイヤを充足させるためだけに人物を補わない．
物語行為$e$が対象化する世界との関係$\mathrm{targ}\subseteq\ko{astav}\times W$は一般に多対多とし，内容の充足からだけでは導出しない（同じ表現を報告として読むか虚構として読むかで対象世界が変わる．メタフィクション理論のtargに当たる）．
作中作は，外側の内容が内側の叙述行為または表現について述べ（aboutness），内側の表現がさらに内容を持つという再帰（\ref{sec:rep-recursion}節のチャネル(a)(c)）で表す．深さはこの経路について数えられるが，存在者に固有のレベル番号としては与えない．外側で「人物が夢を語った」が成立しても，夢の中身を外側へ持ち上げない．

\paragraph{作品} 作品は，制作・公表・継承の歴史によって個体化され，複数の表現を一つの制作単位として束ねる\ko{kika}のタイプ個体である（2026-09-06採択，暫定．devboxの未決を閉じた）．作品は\ko{kilem}と区別し，固有の形態への関数的射影を要求しない（翻訳は形態$F$が違うので，同じ\ko{kilem}の\ko{saden}実現にはならない．これが分離の理由である）．作品と表現の関係は\ko{litoi}ではなく別の関係で与える：
\[
\text{ExpressionOf}\subseteq\ko{kilem}\times\text{作品},\qquad \text{作品実現}(o,x)\ :\Longleftrightarrow\ \exists r\,[\,\text{ExpressionOf}(r,o)\land\ko{saden}(r,x)\,].
\]
\ko{litoi}の右項は集合化された外延なので，作品個体を「その表現の集合」と同一視することになり，歴史的同一性が外延性へ還元されてしまう．作品の同一性は，ある時点で対応する表現の集合の一致によっては決まらない．LRMoo（GOLEMの土台）のWork／Expression／Manifestation／Itemにほぼ対応する．
種水準の述定（「1997年に刊行が始まった」「7作からなる」）は作品に留まり，インスタンス水準の述定（「英語で書かれている」「約7.7万語」）は表現へ分配される．翻訳・改訂を同じ作品に属させる範囲は作品の継承基準によって定め（翻訳／翻案という呼称だけから同一性を決めると循環する），翻案によって別作品を立てる場合には翻案由来$\subseteq$作品$\times$作品を与える．シリーズも，独立した企画・刊行・継承の単位として数える場合に\ko{kika}の個体とし，構成作品への順位付きの関係を与える（「『死の秘宝』は7作目」）．保持（DVD）・提示（上映）の区別は\ko{saden}の制限クラスとして記述し，\ko{takto}／\ko{astav}による網羅的二分とはしない．

\paragraph{世界に依存する分類} \ko{skon}は集合包含であり，その意味と推移律は世界によって変わらない．世界$w$での「$A$ \ko{skon} $B$」は$\mathrm{Ext}_w(A)\subseteq\mathrm{Ext}_w(B)$のことで，外延はフレームの語彙割り当て$\rho_w$が与える（上位分類の章「世界や形式対象」）．作中の語と現実の語が同じ語形を持つことは同じ概念を表す十分条件ではない：同一性条件と定義条件を保つ場合には同じ概念の世界ごとの外延として扱い，条件を変更する場合には別概念を立てる．\ko{coto}は生物学的な人（ホモ・サピエンスという種．\ko{halov}への包摂を定義条件に含む）に固定し，人格・社会的地位による広義の「人」は別語で立てる（未造語．2026-09-07裁定）．したがってマインクラフトのスティーブを「人」と呼ぶことは広義の語の使用であり，\ko{coto}への所属（したがって真核生物への包摂）は導かない．

\paragraph{ゲーム内オブジェクト・理想化された系} ゲーム内オブジェクトは，規則によって計数・持続を定められる世界内個体，作品の設定要素，その状態を表す形式構造，実行時の担体・段階・過程を区別する．同じ世界内個体を保持しながら実行状態が変化し，同じ形式状態が異なる実行で反復されうる．実行過程は\ko{astav}，抽象的な状態値・遷移構造は\ko{stcilkant}，内在的な状態値を持つ個別段階は\ko{canon}，関係的な個別段階は既存関係の引き下げハブである．設定タイプと世界内個体の対応を\ko{saden}で代用しない．
理想化された系（体積$V$の箱に閉じ込められたフェルミ気体）は，それを指定する内容，内容の数学的形式化（形式化$\subseteq$\ko{klaleme}$\times$\ko{stcilkant}，一意でない），仮定の下で対象となる系（箱は\ko{takto}，気体は\ko{matonakt}）を区別する．仮定を満たす世界のクラスは，系そのものでも数学モデルそのものでもない．箱の体積は箱を担い手とする\ko{canon}と量$V$の関係で表し，設定内容に直接\ko{canon}を担わせない．

\begin{example}[ハリー・ポッター]
シリーズ$S$，小説作品$A_1$『賢者の石』，映画作品$A_1^f$（翻案由来$(A_1^f,A_1)$）は作品．英語原文$R_{\mathrm{en}}$・松岡訳$R_{\mathrm{ja}}$は$\text{ExpressionOf}(\cdot,A_1)$，映画本編$R_{\mathrm{film}}$は$\text{ExpressionOf}(\cdot,A_1^f)$．版・冊子・上映・DVDは\ko{saden}．役$T_H$〈ハリー〉はタイプ個体で，対応ハブ$h_{\mathrm{novel}}$（根拠内容$d_1$）と$h_{\mathrm{film}}$（根拠内容$d_1^f$）を持ち，対象はそれぞれ$W_{d_1}$，$W_{d_1^f}$の世界内個体（\ko{takto}$\cap$\ko{halov}$\cap$\ko{kaeno}）．「映画では目が青い」は二つのハブの作中帰結の違いである．設定〈ロンドン〉は現実指示を持ち，$T_H$は持たない（非実在）．ラドクリフは対応者ではなく，演技（\ko{astav}）を介して役に関係する．「ハリー・ポッター」の語はシリーズ・小説作品・原文・役・世界内人物の五義を持ち，持続物の章の見做しグラフで結ぶ．詳細は\texttt{2026-09-06/harry-potter-example.md}．
\end{example}

\begin{whynot}
（2026-09-06）虚構物を\ko{kika}の一種とする案：虚構として記述される対象は人物，出来事，性質，表現，数学的対象と既存分類を横断し，共通の構成射影や同一性基準を供給しない．
GOLEMのクラス階層（Character $\subseteq$ PropositionalObject $\cap$ AgentiveSocialObject 等）を移植する案：物語上の機能や解釈上の区別の多くは既存の個体・二項関係・ロール・語彙仕様で書け，社会的対象としてのキャラクターに行為者的性質を載せると層Kと層Wを一個体に潰して木が混線する．
設定内容，DnSのSituation，正確検証者，充足世界クラスを同一視する案：記述の歴史的同一性，状況の部分構造，正確な成立への寄与，世界クラスの外延的一致は異なる基準である．
全対応者への無条件の述語分配，翻案・類似・対応者関係を同一性として扱う案，作中作のレベルを新しい評価指標とする案，空の充足世界クラスから作中帰結を導く案，対応ハブの根拠を内容から物語行為へ置換する案（行為だけでは内容上の根拠が失われ，一つの行為でも夢・作中作・解釈で対象世界が複数になる），Lewis型の選択関数をフレームに新設する案（背景を一意に選ぶ責任は選択関数の中身に移るだけで消えない），表現\ko{litoi}作品と書く案，保持＝\ko{takto}／提示＝\ko{astav}の網羅的二分，ゲーム内個体を計算状態と同一視する案，理想化系を形式構造と一律に同一視する案，嘘・エラーの専用原始種：いずれも採らない（\texttt{2026-09-06/fiction-codex-log.md}）．
\end{whynot}

\paragraph{語彙テスト用の定義文} 嘘をつく行為は主張行為（\ko{astav}）の定義クラスで，話者が真でないと信じる内容を，相手に真として受け入れさせる意図をもって主張する行為である．内容が実際に偽であることは条件にしない（偽だと信じて述べた内容が偶然真だった場合を落とさないため）．表現・内容そのものに使用機会を離れた嘘の身分を与えない．コンパイルエラーは，検出の発生（\ko{astav}），規則に対するソースの不適合（\ko{kilem}$\times$\ko{klaleme}上の判定関係で，検出を必要条件としない），それを報告する診断表現（\ko{kilem}）の三語義に分ける．診断の出力は不適合の存在をそれだけで含意しない．

\section{抽象・準抽象・具体の境界}
\ko{stcilkant}と\ko{kika}は別物として作る．「この俳句」は\ko{kika}，「この俳句の文字列」は\ko{stcilkant}である．
文字列はずっとあるものだが，俳句はある時点で生み出された．同様に，数式の文字列は\ko{stcilkant}，論文は\ko{kika}，
指示される数学的対象は\ko{stcilkant}である．表現$R$は\ko{stcilkant}の部品$F_{\mathrm{form}}$から作られた\ko{kika}である．

時間的位置だけではチェス，言語，作品，命題，数を分けられないので，次の軸を独立に判定する：
空間的位置（固有・構成・担体依存），歴史的個体化，実現依存（媒体・心的状態・社会的実践），形式的自律性，反復可能性．
\ko{kika}の作業定義（前節の凡例）はこれらのうち歴史的個体化・一般依存・非形式的自律性の組合せである．
命題，アルゴリズム，法律，貨幣，ソフトウェアは反例候補として語彙テストで扱う．

\section{判定テストと典型表現}
定義の明確化のため，言語上の判定テストを置く．
\begin{itemize}
    \item \emph{制作テスト}：「$X$が作られた日付を問えるか」．問える$\to$\ko{kika}．「発見された」が自然$\to$\ko{stcilkant}．
    \item \emph{破損テスト}：「$X$を燃やせる・破けるか」．できる$\to$実現物．
    \item \emph{絶版テスト}：「$X$は絶版になっても存在するか」．する$\to$表現または作品．
    \item \emph{計数テスト}：「同じ$X$が3冊」は実現物を数えている．「$X$の第2版」は表現・作品側．
    \item \emph{翻訳テスト}：「翻訳しても同じ$X$か」．同じ$\to$内容（または作品）．違う$\to$表現．
    \item \emph{偶然同形テスト}：同じ文字列が別言語で別の意味を持つとき，形態は一つ，表現は二つ．
    \item \emph{「では」テスト}：「モールス符号\emph{では}SOSは…」の「では」句が$E_{\mathrm{code}}$の引数位置．鍵・辞書・文法を知らないと読めないとき，参照されているのは$E_{\mathrm{code}}$．
    \item \emph{「作中では」テスト}：付けて真偽が保存される$\to$層W述定．奇妙になる（「作中では1887年に作られた」）$\to$層K述定．
    \item \emph{種水準テスト}：「絶滅する・創作される・翻案される」が当たる$\to$層K（タイプ個体）．
    \item \emph{「として」テスト}：「この岩層は記録として読める」は記述相対の実現．
\end{itemize}
典型的な表現と使う概念の対応：
\begin{itemize}
    \item 「この本は面白い」：実現物の表層から表現・内容への換喩的繰り上げ（善意的解釈）．
    \item 「同じ本を持っている」：表現は同一，実現物は別．
    \item 「彼の発言を引用する」：チャネル(a)．内容に表現が入る．
    \item 「ファイルをUSBへコピーする」：\ko{saden}の連鎖（推移的）．
    \item 「『雪』という語は2文字だ」：形態$F_{\mathrm{form}}$（\ko{stcilkant}）の形式的性質．
    \item 「その詩は1922年に書かれた」：表現$R$（\ko{kika}）への歴史述定．
    \item 「発話」「思考」「演奏」：生起物・実現物・表現・内容を語義分割する．思考の心的プロセスは\ko{anemin}または\ko{astav}，内容だけを\ko{klaleme}とする．
\end{itemize}

\section{語彙テスト}
約200語（言語・記号，作品・芸術，規則・制度，情報・技術，形式対象との境界，心的・社会的境界を均等に）を，
語義と判定用例文の組を単位にボトムアップで検証する．分類表のフィールドは，対象層，空間的位置，歴史的個体化，実現依存，
形式的自律性，反復可能性，形態，内容，符号化規約，aboutness，同一性基準，再帰チャネル，現行候補，判定不能理由とする．
最初の40語をパイロットとし，特徴値を固定した後に残りへ進む．

\begin{devbox}
\textbf{未決事項}
\begin{itemize}
    \item 内容$C_{\mathrm{cont}}$の下位分類（記述的・規則的・設計・物語）と，形式的に自律した内容を\ko{stcilkant}へ分けるか．
    \item 作品は\ko{kilem}と別のタイプ個体として立てた（2026-09-06，暫定）．残るのは作品の継承基準（翻訳・改訂をどこまで同じ作品とするか）．
    \item \ko{kika}を\ko{canon}の担い手として開放するか（作品の語数・指定テンポを値込み\ko{canon}にするか，記述関係のままか）．量・性質モジュールの章は暫定的に「立てない」としており，本章もそれに揃えて未決のまま置く（2026-09-03）．
    \item \ko{saden}の値域に\ko{astav}を含めるか．
    \item 対応ハブ（役・対象・根拠内容・物語行為・解釈行為），背景の持ち込み方針（$U_J$／$\mathrm{Max}$）は暫定．背景候補$\Delta$と優先度の選定基準，\ko{coto}は生物学的種に固定した（2026-09-07）．他の基幹語について何が定義的で何が現実の分類データかは，辞書の\texttt{contents}に定義条件を書く作業として残る．
    \item 未造語：符号化規約，規約参照，aboutness，作品，ExpressionOf，キャラクター役，対応ハブとその射影，登場，翻案由来，現実指示，targ，演じる．
    \item 形式意味論との同定表，YAMATO表現論文との照合，Hjelmslev・Ecoの引用．
\end{itemize}
\end{devbox}

\end{document}
